\naslov{Parametrično reševanje}

Naj bo NDE prvega reda podana implicitno,
\[
  F(x, y, y') = 0.
\]
Denimo, da $y'$ ne moramo eksplicitno izraziti z $x, y$, ali pa je ekspliciten
izraz nepripraven.
Na $F$ poglejmo nekoliko drugače; vsaka dovolj lepa funkcija treh spremenljivk
podaja družino ploskev, $F(\xi, \eta, \zeta) = C$ je implicitna enačba ploskve v
$\R^3$ za vsak $C \in \R$.
Tako podano ploskev lahko parametriziramo.
Naj bo $(u, v) \mapsto (\varphi(u,v), \psi(u,v), \chi(u,v))$ neka
parametrizacija.
Imamo tri pristope za reševanje, v odvisnosti od $F$.

Če $y$ ne nastopa eksplicitno, torej $F(x, y') = 0$, nam enačba definira
krivuljo.
Parametriziramo jo z $\xi = \varphi(t)$ in $\eta = \psi(t)$, da velja
$F(\varphi, \psi) = 0$.
Za poljubno rešitev $t \mapsto (x(t), y(t))$ velja $\dot{y} = y' \dot{x}$, torej
za $\varphi(t) = x(t)$ in $\psi(t) = y'(t)$ dobimo $\dot{y} = \chi
\dot{\varphi}$, oziroma
\[
  y(t) = \int_0^t \chi(\tau) \dot{\varphi}(\tau) d\tau.
\]
Dobimo parametrično izraženo rešitev $t \mapsto (\varphi(t), y(t))$.

\vprasanje{Kako parametrično rešiš enačbo $F(x, y') = 0$?}

Če $x$ ne nastopa eksplicitno, torej $F(y, y') = 0$, dobimo enačbo krivulje
$F(\xi, \eta) = 0$, ki jo parametriziramo s $t \mapsto (\chi(t), \psi(t))$.
Če označimo $\psi = y$ in $\chi = y'$, velja $\dot{\psi} = \chi \dot{x}$ oziroma
\[
  x(t) = \int_0^t \frac{\dot{\psi}(\tau)}{\chi(\tau)} d\tau,
\]
torej je $t \mapsto (x(t), \psi(t))$ parametrično podana rešitev.

\vprasanje{Kako parametrično rešiš enačbo $F(y, y') = 0$?}

V splošnem nam $F(x, y, y') = 0$ definira ploskev.
Parametriziramo jo kot zgoraj z
\begin{align*}
  x = \varphi(u, v) && y = \psi(u, v) && y' = \chi(u,v)
\end{align*}
Naj bo $t \mapsto (x(u(t), v(t)), y(u(t), v(t)))$ neka rešitev naše enačbe.
Potem je $t \mapsto (\varphi, \psi, \chi)$ krivulja na ploskvi.
V nadaljevanju predpostavimo, da je preslikava $(u,v) \mapsto (x, y)$ obrnljiva,
in izračunajmo
\begin{gather*}
  \dot{y} = y_u \dot{u} + y_v \dot{v} = \psi_u \dot{u} + \psi_v \dot{v}, \\
  \dot{x} = \varphi_u \dot{u} + \varphi_v \dot{v}.
\end{gather*}
Ker tudi tu velja $\dot{y} = y' \dot{x}$, izrazimo
\[
  u' = \frac{\dot{u}}{\dot{v}} = - \frac{\psi_v - \chi \varphi_v}{\psi_u - \chi
	\varphi_u}.
\]
Dobili smo eksplicitno enačbo prvega reda v spremenljivki $u = u(v)$.

\vprasanje{Kako parametrično rešiš $F(x, y, y') = 0$?}

\podnaslov{Lagrangeova in Clairontova enačba}

Lagrangeova enačba je enačba oblike
\[
  y = x \varphi(y') + \psi(y').
\]
Rešujemo jo parametrično; $x = u$, $y' = v$ in $y = u \varphi(v) + \psi(v)$.
Velja $dy = y' dx$, iz česar izpeljemo
\[
  (\varphi(v) - v) du + (u \varphi'(v) + \psi'(v)) dv = 0.
\]
Če je $\varphi(v) \ne v$, dobimo
\[
  (\phi(v) - v) \frac{du}{dv} + u \varphi'(v) + \psi'(v) = 0,
\]
kar je linearna diferencialna enačba prvega reda, če pa je $\varphi(v) = v$, pa
imamo Clairontovo enačbo
\[
  y = x y' + \psi(y').
\]
To predelamo v
\[
  (u + \psi'(v)) dv = 0,
\]
in obravnavamo dva primera.
Če je $dv = 0$, je $y'$ konstanta, torej dobimo družino rešitev $y = Cx +
\psi(C)$ (to vstavimo v enačbo; ni nujno vsak $C$ dober).
Če pa je $u + \psi'(v) = 0$, pa dobimo še eno rešitev.

\vprasanje{Kaj sta Lagrangeova in Clairontova enačba? Kako ju rešimo?}

\podnaslov{Ovojnice družin krivulj}

Imejmo družino krivulj, podano implicitno z enačbo $F(x, y, C) = 0$.
Denimo, da je družina taka, da obstaja krivulja, ki se v vsaki svoji točki
dotika natanko enega člana družine.
Taki krivulji pravimo \pojem{ovojnica} družine.
Smiselno jo je parametrizirati z $C \mapsto (x(C), y(C))$, pri čemer se ovojnica
v točki $(x(C), y(C))$ dotika člana družine s tem $C$.
Denimo, da je parametrizacija regularna, torej za vsak $C$
\[
  (\partial_C x)^2 + (\partial_C y)^2 \ne 0.
\]
Definirajmo
\[
  \phi(C) = F(x(C), y(C), C).
\]
Ker funkcija izračuna $F$ v točki na krivulji, je $\phi = 0$.
Torej
\[
  \phi'(C) = \partial_x F \partial_C x + \partial_y F \partial_C y + \partial_C
  F = 0.
\]

Če je $t \mapsto (x(t), y(t))$ parametrizacija $C_0$-tega člana družine, velja
\[
  \partial_t F(x(t), y(t), C_0)
  = \partial_x F \dot{x} + \partial_y F \dot{y} = 0
\]
v točki dotika z ovojnico.
Vektor $[\dot{x}, \dot{y}]^T$ je vzporeden z $[\partial_C x, \partial_C y]^T$ v
tej točki, torej je
\[
  \begin{bmatrix}
	\partial_x F \\ \partial_y F
  \end{bmatrix}
  \bot
  \begin{bmatrix}
	\dot{x} \\ \dot{y}
  \end{bmatrix}
  \parallel
  \begin{bmatrix}
	\partial_C x \\ \partial_C y
  \end{bmatrix}
\]
in zato
\[
  \partial_x F \partial_C x + \partial_y F \partial_C y = 0.
\]
Torej v točki dotika velja
\[
  \partial_C F = 0.
\]
Iz para enačb $F(x, y, C) = 0$ in $\partial_C F = 0$ dobimo vse točke na
ovojnici.

\vprasanje{Kaj je ovojnica družine krivulj? Kako jo izračunaš? Izpelji.}

% LocalWords:  nepriraven
